\section{Introduction}
\label{sec:introduction}

Procedural memories and reusable agent skills can improve software work, but a
procedure that was safe in one context may be unsafe when a relevant
precondition changes.  The causal question is sharper than whether a memory is
topically relevant or whether an endpoint reports success: does assigning a
demonstrably source-correct and source-focal-safe procedure change
task-complete insecurity when its security-relevant applicability condition
does not transfer?

We initially designed V1 to estimate the causal contrast. A systematic audit
showed that it did not identify that quantity; we therefore use it here as a
formative case study of evaluation misidentification.  Its nominal
functionality endpoint accepted unchanged repositories, memory delivery
displaced usable trajectory resources, focal security was at floor, preserved
traces showed no substantive procedural uptake, and nominal seed repetitions
were nearly deterministic.  These observations do not estimate a harmful
memory effect.  They expose how a plausible evaluation can answer a different
question from the one its design intends.

We then made the evidentiary implications of the intended estimand explicit.
The resulting contract ties target nontriviality, source safety, directional
applicability, treatment isolation, process measurement, and joint
task-security outcomes to the causal statements each licenses.  A bounded
real-repository construction study applied these gates prospectively and
stopped without an admissible complete pair.  A deliberately constructed
positive control subsequently traversed the unchanged path, while two
single-gate near misses were rejected at their intended gates.  Finally, a
frozen six-work external framework application found the combined framework
partially distinctive while documenting substantial antecedents and
legitimately different estimands in prior work.

\paragraph{Contributions.}
The paper provides four methodological contributions:
\begin{enumerate}[leftmargin=*,itemsep=2pt,topsep=3pt]
  \item an estimand-derived identification contract for security-sensitive
  procedural-memory transfer;
  \item a formative evaluation audit showing that nominal endpoint success can
  coexist with no identified task completion, no substantive observable
  memory uptake, resource displacement, a security-floor effect, and
  pseudo-replication;
  \item a prospectively bounded, artifact-grounded real-repository construction
  study that records gate-specific attrition and an explicit stop decision;
  and
  \item non-vacuity and external validation through a known-valid control and a
  six-paper framework application.
\end{enumerate}
The artifact machinery is part of these methodological claims: hashes, source
locks, executable controls, content-access records, and stopping ledgers bind
each inferential requirement to inspectable evidence.
